% GENERATED FILE. Edit canonical lessons, then run npm run generate:books.
% canonical-source-hash: fnv1a64:8d9b8dbd7876933a
% canonical-lessons: MW-C171-basan, MW-C171-chithro, MW-C171-majno, MW-C171-saph, MW-C171-kachro

\chapter{Washing Up}
\label{ch:mw-washing-up}
\hlchaptermodality{171}

\begin{chapteropening}
\textbf{By the end of this chapter:} \emph{I can say how I clean the pots and the kitchen.}

Complete the last lesson of chapter 171: I can say how I clean the pots and the kitchen.
\end{chapteropening}

\section[\texorpdfstring{\mw{बासण} (bāsaṇ)}{बासण (bāsaṇ)}]{\mw{बासण} (bāsaṇ) — a utensil; a pot}

\label{lesson:MW-C171-basan}



\begin{quote}
Before the new one: say the Marwadi for to dry (clothes), then the Marwadi for to gather up.
\end{quote}

\subsection*{You'll want to know: \mw{बासण}}
\textbf{\mw{बासण}} — \emph{bāsaṇ} — \textquotedblleft{}a utensil; a pot\textquotedblright{}.

\begin{quote}
\mw{रसोई} \mw{मांय} \mw{घणा} \mw{बासण} \mw{है।}
\end{quote}

\begin{quote}
\emph{rasoī mā̃y ghaṇā bāsaṇ hai.}
\end{quote}

\begin{quote}
There are a lot of pots in the kitchen.
\end{quote}

A masculine noun.

\begin{cousinweb}
\textbf{\mw{बासण}} (\emph{bāsaṇ}) is inherited from Sanskrit \emph{vāsana}, \textquotedblleft{}a box, a casket\textquotedblright{}; CDIAL lists Hindi \emph{bāsan} and Gujarati \emph{vāsaṇ}, \textquotedblleft{}a pot, a dish\textquotedblright{} (CDIAL 11599). CDIAL gives no Marwari form in this entry.
\end{cousinweb}

\begin{grammarlens}[title={What you've built}]
One more word for saying what you do.
\end{grammarlens}

\subsection*{Your turn}
\begin{itemize}
\raggedright
  \item \emph{Say it:} \emph{bāsaṇ}
  \item \emph{Say it:} \emph{bāsaṇ}, once more
  \item \emph{Say it:} \emph{rasoī mā̃y ghaṇā bāsaṇ hai.} — There are a lot of pots in the kitchen.
  \item \emph{Recall:} say the Marwadi for to dry (clothes), then the Marwadi for to gather up, then say \emph{bāsaṇ} again
\end{itemize}

\begin{tcolorbox}[breakable,colback=teal!4,colframe=teal!35!black,title={Before you move on}]
What does \mw{बासण} mean? (\textquotedblleft{}A utensil; a pot\textquotedblright{}.) Say it once more.
\end{tcolorbox}
\section[\texorpdfstring{\mw{चीथड़ो} (chīthṛo)}{चीथड़ो (chīthṛo)}]{\mw{चीथड़ो} (chīthṛo) — a rag}

\label{lesson:MW-C171-chithro}



\begin{quote}
Before the new one: say the Marwadi for to gather up, then the Marwadi for a utensil.
\end{quote}

\subsection*{You'll want to know: \mw{चीथड़ो}}
\textbf{\mw{चीथड़ो}} — \emph{chīthṛo} — \textquotedblleft{}a rag\textquotedblright{}.

\begin{quote}
\mw{चीथड़ो} \mw{मैलो} \mw{है।}
\end{quote}

\begin{quote}
\emph{chīthṛo mailo hai.}
\end{quote}

\begin{quote}
The rag is dirty.
\end{quote}

A masculine noun.

\begin{cousinweb}
\textbf{\mw{चीथड़ो}} (\emph{chīthṛo}) comes from a reconstructed \emph{*citth}, \textquotedblleft{}pound, tear\textquotedblright{}; CDIAL lists Hindi \emph{cīthṛā} and Sindhi \emph{cīthaṛo}, both masculine, for a rag or old torn clothes (CDIAL 4802). CDIAL gives no Marwari form in this entry.
\end{cousinweb}

\begin{grammarlens}[title={What you've built}]
One more word for saying what you do.
\end{grammarlens}

\subsection*{Your turn}
\begin{itemize}
\raggedright
  \item \emph{Say it:} \emph{chīthṛo}
  \item \emph{Say it:} \emph{chīthṛo}, once more
  \item \emph{Say it:} \emph{chīthṛo mailo hai.} — The rag is dirty.
  \item \emph{Recall:} say the Marwadi for to gather up, then the Marwadi for a utensil, then say \emph{chīthṛo} again
\end{itemize}

\begin{tcolorbox}[breakable,colback=teal!4,colframe=teal!35!black,title={Before you move on}]
What does \mw{चीथड़ो} mean? (\textquotedblleft{}A rag\textquotedblright{}.) Say it once more.
\end{tcolorbox}
\section[\texorpdfstring{\mw{मांजणो} (mā̃jṇo)}{मांजणो (mā̃jṇo)}]{\mw{मांजणो} (mā̃jṇo) — to scrub (pots)}

\label{lesson:MW-C171-majno}



\begin{quote}
Before the new one: say the Marwadi for a utensil, then the Marwadi for a rag.
\end{quote}

\subsection*{You'll want to know: \mw{मांजणो}}
\textbf{\mw{मांजणो}} — \emph{mā̃jṇo} — \textquotedblleft{}to scrub (pots)\textquotedblright{}.

\begin{quote}
\mw{हूं} \mw{आथण} \mw{बासण} \mw{मांजूं} \mw{हूं।}
\end{quote}

\begin{quote}
\emph{hū̃ āthaṇ bāsaṇ mā̃jū̃ hū̃.}
\end{quote}

\begin{quote}
I scrub the pots in the evening.
\end{quote}

Past (I …): \mw{म्हैं} \mw{बासण} \mw{मांज्या}, \emph{mhaĩ bāsaṇ mā̃jyā}.

\begin{cousinweb}
\textbf{\mw{मांजणो}} (\emph{mā̃jṇo}) is inherited from Sanskrit \emph{mārjati}, \textquotedblleft{}rubs clean, wipes off\textquotedblright{}; CDIAL lists Hindi \emph{mā̃jnā} and Gujarati \emph{mā̃jvũ} (CDIAL 10080). CDIAL gives no Marwari form in this entry.
\end{cousinweb}

\begin{grammarlens}[title={What you've built}]
One more word for saying what you do.
\end{grammarlens}

\subsection*{Your turn}
\begin{itemize}
\raggedright
  \item \emph{Say it:} \emph{mā̃jṇo}
  \item \emph{Say it:} \emph{mā̃jṇo}, once more
  \item \emph{Say it:} \emph{hū̃ āthaṇ bāsaṇ mā̃jū̃ hū̃.} — I scrub the pots in the evening.
  \item \emph{Recall:} say the Marwadi for a utensil, then the Marwadi for a rag, then say \emph{mā̃jṇo} again
\end{itemize}

\begin{tcolorbox}[breakable,colback=teal!4,colframe=teal!35!black,title={Before you move on}]
What does \mw{मांजणो} mean? (\textquotedblleft{}To scrub (pots)\textquotedblright{}.) Say it once more.
\end{tcolorbox}
\section[\texorpdfstring{\mw{साफ} (sāph)}{साफ (sāph)}]{\mw{साफ} (sāph) — clean}

\label{lesson:MW-C171-saph}



\begin{quote}
Before the new one: say the Marwadi for a rag, then the Marwadi for to scrub.
\end{quote}

\subsection*{You'll want to know: \mw{साफ}}
\textbf{\mw{साफ}} — \emph{sāph} — \textquotedblleft{}clean\textquotedblright{}.

\begin{quote}
\mw{म्हारी} \mw{थाळी} \mw{साफ} \mw{है।}
\end{quote}

\begin{quote}
\emph{mhārī thāḷī sāph hai.}
\end{quote}

\begin{quote}
My plate is clean.
\end{quote}

\begin{cousinweb}
\textbf{\mw{साफ}} (\emph{sāph}) is a loan, usually traced to Persian \emph{sāf}, \textquotedblleft{}clean\textquotedblright{}, which Persian took from Arabic \emph{ṣāfī}, \textquotedblleft{}pure, clear\textquotedblright{}. It does not change: \textbf{\mw{साफ} \mw{कपड़ा}}, \textbf{\mw{साफ} \mw{थाळी}}.
\end{cousinweb}

\begin{grammarlens}[title={What you've built}]
One more word for saying what you do.
\end{grammarlens}

\subsection*{Your turn}
\begin{itemize}
\raggedright
  \item \emph{Say it:} \emph{sāph}
  \item \emph{Say it:} \emph{sāph}, once more
  \item \emph{Say it:} \emph{mhārī thāḷī sāph hai.} — My plate is clean.
  \item \emph{Recall:} say the Marwadi for a rag, then the Marwadi for to scrub, then say \emph{sāph} again
\end{itemize}

\begin{tcolorbox}[breakable,colback=teal!4,colframe=teal!35!black,title={Before you move on}]
What does \mw{साफ} mean? (\textquotedblleft{}Clean\textquotedblright{}.) Say it once more.
\end{tcolorbox}
\section[\texorpdfstring{\mw{कचरो} (kachro)}{कचरो (kachro)}]{\mw{कचरो} (kachro) — rubbish}

\label{lesson:MW-C171-kachro}



\begin{quote}
Before the new one: say the Marwadi for to scrub, then the Marwadi for clean.
\end{quote}

\subsection*{You'll want to know: \mw{कचरो}}
\textbf{\mw{कचरो}} — \emph{kachro} — \textquotedblleft{}rubbish\textquotedblright{}.

\begin{quote}
\mw{रसोई} \mw{मांय} \mw{कचरो} \mw{कोनी।}
\end{quote}

\begin{quote}
\emph{rasoī mā̃y kachro konī.}
\end{quote}

\begin{quote}
There's no rubbish in the kitchen.
\end{quote}

A masculine noun.

\begin{cousinweb}
\textbf{\mw{कचरो}} (\emph{kachro}) is inherited from Sanskrit \emph{kaccara}, \textquotedblleft{}dirty\textquotedblright{}; CDIAL lists Hindi \emph{kacrā}, \textquotedblleft{}sweepings, dirt\textquotedblright{} (CDIAL 2615). CDIAL gives no Marwari form in this entry.
\end{cousinweb}

\begin{grammarlens}[title={What you've built}]
One more word for saying what you do.
\end{grammarlens}

\subsection*{Your turn}
\begin{itemize}
\raggedright
  \item \emph{Say it:} \emph{kachro}
  \item \emph{Say it:} \emph{kachro}, once more
  \item \emph{Say it:} \emph{rasoī mā̃y kachro konī.} — There's no rubbish in the kitchen.
  \item \emph{Recall:} say the Marwadi for to scrub, then the Marwadi for clean, then say \emph{kachro} again
\end{itemize}

\begin{tcolorbox}[breakable,colback=teal!4,colframe=teal!35!black,title={Before you move on}]
What does \mw{कचरो} mean? (\textquotedblleft{}Rubbish\textquotedblright{}.) Say it once more.
\end{tcolorbox}
